\section{Declaración de uso de inteligencia artificial}
\label{sec:declaracion-ia}

{\small
\begin{xltabular}{\textwidth}{@{}>{\raggedright\arraybackslash}p{2.5cm}>{\raggedright\arraybackslash}p{1.7cm}LL@{}}
  \caption{Herramientas de inteligencia artificial usadas como apoyo}
  \label{tab:declaracion-ia} \\
  \toprule
  \textbf{Herramienta} & \textbf{Versión} & \textbf{Para qué se usó} & \textbf{Cómo se verificó lo obtenido} \\
  \midrule
  \endfirsthead
  \toprule
  \textbf{Herramienta} & \textbf{Versión} & \textbf{Para qué se usó} & \textbf{Cómo se verificó lo obtenido} \\
  \midrule
  \endhead
  \bottomrule
  \endlastfoot
    ChatGPT (OpenAI), Steffanny Loor & GPT-6 & Apoyo para organizar el árbol de fallos, mejorar la redacción del análisis de causas raíz y revisar
    la presentación del diagrama. & Se contrastó la información con el informe oficial del incidente de GitHub de 2018 \cite{Warner2018}. También se revisó el diagrama y se compiló el documento para comprobar su presentación. \\
    \addlinespace
    Claude Code (Anthropic), aplicación de escritorio, Mario Castro & Claude Fable 5.2 (primer paso) y Claude Opus 5.5 (resto) & Plan de trabajo a partir del enunciado; comparación de los informes oficiales de los cuatro casos y recomendación de uno; carátula del SGA con el script del curso; estructura vacía del repositorio, ramas y configuración de fusiones; mensajes de coordinación del grupo; revisión de los datos del informe de GitHub y su entrada BibLaTeX; entradas de las dos referencias de CAP con \texttt{doi2bib}; revisión de la sección 1, de la fila de lecciones propia y de los PR del equipo contra el informe; conflictos de la tabla de lecciones; ajuste de formato para que el análisis quepa en 2 páginas. & Los comandos se ejecutaron manualmente y el resultado se comprobó en GitHub. Cada observación de las revisiones se comparó con el informe. Las referencias con DOI se contrastaron con doi.org. El documento se compiló con los comandos del README. \\
    \addlinespace
    ChatGPT (OpenAI), Luis Quintero & GPT-6 & Organizar las secciones 3 y 4, elaborar un borrador inicial y revisar la redacción y la precisión técnica del análisis. También ayudó a localizar las cifras del informe de GitHub y las referencias académicas sobre CAP. & Las cifras se contrastaron con el informe oficial, en especial con la cronología del 22 de octubre a las 16:45 UTC. La explicación de CAP se revisó con las partes del informe sobre la partición, la recuperación de las bases de datos y las réplicas atrasadas, y con los artículos de Gilbert y Lynch (2002) y Brewer (2012). La versión final se revisó para conservar solo la información necesaria y las referencias que respaldan el análisis. \\
    \addlinespace
    Claude (Anthropic), Dayver Rivas & Claude Opus 5.5 & Sección 5: leyó el informe oficial, sacó los puntos de mitigación y redactó un primer borrador, que se subió al PR; después Dayver reescribió la sección con sus palabras y Claude revisó sus versiones contra el informe (datos inexactos, extensión y uso de \texttt{\textbackslash ref}). Fila de lecciones: revisó el repositorio de su PFC, propuso dos riesgos posibles (Dayver eligió uno), agregó la plantilla vacía de la fila y revisó su redacción. & Comparó cada dato de la sección 5 con el informe y corrigió lo que no coincidía, por ejemplo «conciliarlas manualmente» y «se reconfiguró Orchestrator». Para la fila de lecciones, comprobó en el código de su PFC que la subida del certificado y el registro en la base se hacen en dos pasos separados. \\
\end{xltabular}
}

\noindent\textbf{Ejemplo de prompt (Steffanny Loor):} ``Ayúdame a revisar las causas raíz del incidente de GitHub de octubre de 2018, tomando como referencia el informe oficial. Quiero organizar las causas en un árbol de fallos, identificar los problemas de replicación y conmutación automática, y comprobar que la explicación sea clara y técnicamente correcta.''
